%% chapters/ch2_background.tex

\chapter{文獻探討}
\label{ch:background}

\section{相關研究}
\label{sec:related-work}

這一節回顧與本研究相關的既有工作。多筆文獻可以一次引用，
例如 \cite{knuth1984texbook, mittelbach2004companion, shannon1948mathematical}，
\verb|cite| 套件會把三筆以上的連號合併成區間。

\section{理論基礎}
\label{sec:theory}

這一節整理後續章節會用到的理論。多行推導建議用 \verb|align| 環境，
只在需要交互參考的行加 \verb|\label|：

\begin{align}
  X(f) &= \int_{-\infty}^{\infty} x(t)\, e^{-j 2\pi f t}\, \mathrm{d}t
  \label{eq:fourier} \\
  x(t) &= \int_{-\infty}^{\infty} X(f)\, e^{ j 2\pi f t}\, \mathrm{d}f
  \label{eq:inverse-fourier}
\end{align}

式 \eqref{eq:fourier} 與式 \eqref{eq:inverse-fourier} 互為反轉換；
資訊量的定義見 \cite{shannon1948mathematical}，
離散版本的快速演算法見 \cite{cooley1965algorithm}。

\subsection{符號慣例}
\label{subsec:notation}

全篇使用的符號整理於前頁的符號說明。向量以粗體小寫表示，
矩陣以粗體大寫表示。
